\section{Overall Project Context}\label{sec:overall-project-context}

Modern security operations depend on specialised tooling to discover, assess and
remediate faults across networks and applications. As the scale of this work
grows, the manual translation of an operator's intent into safe, well-scoped tool
actions becomes a significant bottleneck, and mistakes made at this point---an
incorrect target, a malformed instruction, or an out-of-scope action---can affect
production infrastructure. Large Language Models (LLMs) are attractive for this
translation step because they can interpret natural language and emit structured,
machine-readable instructions that drive security tooling
\parencite{deng2024pentestgpt}. Yet a model is a stochastic component: its output
may be malformed, misclassified, hallucinated, or the product of injected
instructions, and a security pipeline that acts on such output without further
checking transfers every one of these errors into the actions it executes.

NeuroShell is an autonomous, natural-language-driven security orchestration
framework developed as a four-component research chain: the Intent Recognition
Engine (IRE), a Dynamic Planner, an Adaptive Execution and Error Recovery engine,
and an AI-driven Vulnerability Analysis engine (project reference R26-CS-018).
The framework is designed to let an analyst express offline security tasks in
natural language and have the resulting validated, structured contracts drive
downstream automation, within an explicitly configured engagement scope. It
therefore sits at the intersection of cybersecurity automation, software
engineering and applied machine learning. Its central concern is the first step
of that chain: the point at which human intent is converted into a
machine-executable form, because an error introduced there propagates through
every stage that follows and must be detected before it takes effect.

This dissertation concentrates on a single component of the framework and asks
whether, in an offline CPU-only deployment, a small instruction-tuned language
model can be made sufficiently reliable to serve as the head of an automated
security pipeline when its raw output is constrained by deterministic
validation.